\documentclass[10pt,a4paper]{amsart}
\usepackage{amsmath,amssymb,amsthm,graphicx}
\usepackage{verbatim}
\newenvironment{reference}{\comment}{\endcomment}
%\newenvironment{reference}{}{}
\newenvironment{slogan}{\comment}{\endcomment}
\newenvironment{history}{\comment}{\endcomment}

% For commutative diagrams we use Xy-pic
\usepackage[all]{xy}

% We use 2cell for 2-commutative diagrams.
\xyoption{2cell}
% Theorem environments.
%
\theoremstyle{plain}
\newtheorem{theorem}[subsection]{Theorem}
\newtheorem{proposition}[subsection]{Proposition}
\newtheorem{lemma}[subsection]{Lemma}

\theoremstyle{definition}
\newtheorem{definition}[subsection]{Definition}
\newtheorem{example}[subsection]{Example}
\newtheorem{exercise}[subsection]{Exercise}
\newtheorem{situation}[subsection]{Situation}

\theoremstyle{remark}
\newtheorem{remark}[subsection]{Remark}
\newtheorem{remarks}[subsection]{Remarks}

\numberwithin{equation}{subsection}

% Macros
%
\def\lim{\mathop{\mathrm{lim}}\nolimits}
\def\colim{\mathop{\mathrm{colim}}\nolimits}
\def\Spec{\mathop{\mathrm{Spec}}}
\def\Hom{\mathop{\mathrm{Hom}}\nolimits}
\def\Ext{\mathop{\mathrm{Ext}}\nolimits}
\def\SheafHom{\mathop{\mathcal{H}\!\mathit{om}}\nolimits}
\def\SheafExt{\mathop{\mathcal{E}\!\mathit{xt}}\nolimits}
\def\Sch{\mathit{Sch}}
\def\Mor{\mathop{\mathrm{Mor}}\nolimits}
\def\Ob{\mathop{\mathrm{Ob}}\nolimits}
\def\Sh{\mathop{\mathit{Sh}}\nolimits}
\def\NL{\mathop{N\!L}\nolimits}
\def\CH{\mathop{\mathrm{CH}}\nolimits}
\def\proetale{{pro\text{-}\acute{e}tale}}
\def\etale{{\acute{e}tale}}
\def\QCoh{\mathit{QCoh}}
\def\Ker{\mathop{\mathrm{Ker}}}
\def\Im{\mathop{\mathrm{Im}}}
\def\Coker{\mathop{\mathrm{Coker}}}
\def\Coim{\mathop{\mathrm{Coim}}}

% Boxtimes
%
\DeclareMathSymbol{\boxtimes}{\mathbin}{AMSa}{"02}

%
% Macros for moduli stacks/spaces
%
\def\QCohstack{\mathcal{QC}\!\mathit{oh}}
\def\Cohstack{\mathcal{C}\!\mathit{oh}}
\def\Spacesstack{\mathcal{S}\!\mathit{paces}}
\def\Quotfunctor{\mathrm{Quot}}
\def\Hilbfunctor{\mathrm{Hilb}}
\def\Curvesstack{\mathcal{C}\!\mathit{urves}}
\def\Polarizedstack{\mathcal{P}\!\mathit{olarized}}
\def\Complexesstack{\mathcal{C}\!\mathit{omplexes}}
% \Pic is the operator that assigns to X its picard group, usage \Pic(X)
% \Picardstack_{X/B} denotes the Picard stack of X over B
% \Picardfunctor_{X/B} denotes the Picard functor of X over B
\def\Pic{\mathop{\mathrm{Pic}}\nolimits}
\def\Picardstack{\mathcal{P}\!\mathit{ic}}
\def\Picardfunctor{\mathrm{Pic}}
\def\Deformationcategory{\mathcal{D}\!\mathit{ef}}
\usepackage{cleveref}
\crefname{originaltheorem}{original theorem}{original theorems}
\Crefname{originaltheorem}{Original theorem}{Original theorems}
\makeatletter
\newlabel{etale-theorem-remarkable-equivalence}{{039R}{source}{}{sourcepointer-1}{}}
\newlabel{etale-theorem-remarkable-equivalence@cref}{{[originaltheorem][0][]039R}{[1][0][]source}{}{}{}}
\makeatother
\crefname{originallemma}{original lemma}{original lemmas}
\Crefname{originallemma}{Original lemma}{Original lemmas}
\makeatletter
\newlabel{more-morphisms-lemma-deform-property}{{06AG}{source}{}{sourcepointer-2}{}}
\newlabel{more-morphisms-lemma-deform-property@cref}{{[originallemma][0][]06AG}{[1][0][]source}{}{}{}}
\makeatother
\crefname{originalproposition}{original proposition}{original propositions}
\Crefname{originalproposition}{Original proposition}{Original propositions}
\makeatletter
\newlabel{spaces-limits-proposition-affine}{{07VT}{source}{}{sourcepointer-3}{}}
\newlabel{spaces-limits-proposition-affine@cref}{{[originalproposition][0][]07VT}{[1][0][]source}{}{}{}}
\makeatother
\crefname{originaldefinition}{original definition}{original definitions}
\Crefname{originaldefinition}{Original definition}{Original definitions}
\makeatletter
\newlabel{spaces-perfect-definition-elementary-distinguished-square}{{08GM}{source}{}{sourcepointer-4}{}}
\newlabel{spaces-perfect-definition-elementary-distinguished-square@cref}{{[originaldefinition][0][]08GM}{[1][0][]source}{}{}{}}
\makeatother
\crefname{originallemma}{original lemma}{original lemmas}
\Crefname{originallemma}{Original lemma}{Original lemmas}
\makeatletter
\newlabel{spaces-perfect-lemma-make-more-elementary-distinguished-squares}{{08GN}{source}{}{sourcepointer-5}{}}
\newlabel{spaces-perfect-lemma-make-more-elementary-distinguished-squares@cref}{{[originallemma][0][]08GN}{[1][0][]source}{}{}{}}
\makeatother
\crefname{originallemma}{original lemma}{original lemmas}
\Crefname{originallemma}{Original lemma}{Original lemmas}
\makeatletter
\newlabel{spaces-pushouts-lemma-construct-elementary-distinguished-square}{{0DVJ}{source}{}{sourcepointer-6}{}}
\newlabel{spaces-pushouts-lemma-construct-elementary-distinguished-square@cref}{{[originallemma][0][]0DVJ}{[1][0][]source}{}{}{}}
\makeatother
\crefname{originallemma}{original lemma}{original lemmas}
\Crefname{originallemma}{Original lemma}{Original lemmas}
\makeatletter
\newlabel{spaces-pushouts-lemma-elementary-distinguished-square-pushout}{{0DVI}{source}{}{sourcepointer-7}{}}
\newlabel{spaces-pushouts-lemma-elementary-distinguished-square-pushout@cref}{{[originallemma][0][]0DVI}{[1][0][]source}{}{}{}}
\makeatother
\title{Etale lifts of approximations across nilpotent thickenings with prescribed special fibre}
\author{The Stacks Project Authors}
\date{stacks-project:restricted.tex:local-snapshot-20261012; GNU FDL 1.2 or later}
\begin{document}
\maketitle
\noindent\textit{Editorial note.} All delivered mathematical text is the author's own, copied from the pinned Stacks Project chapter file restricted.tex; no statement or proof step was rewritten, reordered or removed.\par\medskip
\noindent\textit{Editorial note.} The counted proof is the author's complete argument: a distinguished-square recognition on the special fibre, a deformation of the etale leg by topological invariance, a pushout construction pinned by base-change uniqueness, and an affineness conclusion, all inside the delivered span\par\medskip
\noindent\textit{Editorial note.} The label lemma-lift-approximation names the counted statement, and the proof environment immediately following it in the source is delivered verbatim as its proof.\par\medskip
\noindent\textit{Editorial note.} The counted statement is carried by the author's own lemma environment in the chapter source.\par\medskip
\noindent\textit{Editorial note.} The proof closes with the author's sentence 'It is clear that B fits into the desired diagram and satisfies all the properties required of it.'; the properties were established in the course of the construction, and this closing verification is retained verbatim.\par\medskip
\noindent\textit{Allowed standard prerequisites (editorial).}\begin{itemize}
\item Elementary distinguished squares of algebraic spaces and their base change stability
\item Topological invariance of the etale site under finite order thickenings, with separatedness deformed along thickenings
\item The pushout construction of elementary distinguished squares in the pushouts theory, with the identification of the pushout after base change, and affineness transfer from Limits of Spaces
\end{itemize}\medskip
\ifdefined\hypertarget\hypertarget{sourcepointer-1}{}\fi
\noindent\textit{Source reference.} Original theorem 039R: sources/stacks-chow/tags.txt, source lines 4282--4282. Prerequisite pointer invoked at proof line(s) 2219; it names a standard fact of a sibling chapter that the counted argument uses, and every step of the counted proof built on it is retained verbatim in the delivered body.\par\medskip
\ifdefined\hypertarget\hypertarget{sourcepointer-2}{}\fi
\noindent\textit{Source reference.} Original lemma 06AG: sources/stacks-chow/tags.txt, source lines 8075--8075. Prerequisite pointer invoked at proof line(s) 2222; it names a standard fact of a sibling chapter that the counted argument uses, and every step of the counted proof built on it is retained verbatim in the delivered body.\par\medskip
\ifdefined\hypertarget\hypertarget{sourcepointer-3}{}\fi
\noindent\textit{Source reference.} Original proposition 07VT: sources/stacks-chow/tags.txt, source lines 10018--10018. Prerequisite pointer invoked at proof line(s) 2249; it names a standard fact of a sibling chapter that the counted argument uses, and every step of the counted proof built on it is retained verbatim in the delivered body.\par\medskip
\ifdefined\hypertarget\hypertarget{sourcepointer-4}{}\fi
\noindent\textit{Source reference.} Original definition 08GM: sources/stacks-chow/tags.txt, source lines 10750--10750. Prerequisite pointer invoked at proof line(s) 2210; it names a standard fact of a sibling chapter that the counted argument uses, and every step of the counted proof built on it is retained verbatim in the delivered body.\par\medskip
\ifdefined\hypertarget\hypertarget{sourcepointer-5}{}\fi
\noindent\textit{Source reference.} Original lemma 08GN: sources/stacks-chow/tags.txt, source lines 10751--10751. Prerequisite pointer invoked at proof line(s) 2235; it names a standard fact of a sibling chapter that the counted argument uses, and every step of the counted proof built on it is retained verbatim in the delivered body.\par\medskip
\ifdefined\hypertarget\hypertarget{sourcepointer-6}{}\fi
\noindent\textit{Source reference.} Original lemma 0DVJ: sources/stacks-chow/tags.txt, source lines 17364--17364. Prerequisite pointer invoked at proof line(s) 2224; it names a standard fact of a sibling chapter that the counted argument uses, and every step of the counted proof built on it is retained verbatim in the delivered body.\par\medskip
\ifdefined\hypertarget\hypertarget{sourcepointer-7}{}\fi
\noindent\textit{Source reference.} Original lemma 0DVI: sources/stacks-chow/tags.txt, source lines 17363--17363. Prerequisite pointer invoked at proof line(s) 2247; it names a standard fact of a sibling chapter that the counted argument uses, and every step of the counted proof built on it is retained verbatim in the delivered body.\par\medskip
\begin{lemma}
\label{lemma-lift-approximation}
Let $A$ be a Noetherian ring and $I \subset A$ an ideal.
Let $J \subset A$ be a nilpotent ideal. Consider a commutative diagram
$$
\xymatrix{
C \ar[r] & C_0 \ar@{=}[r] & C/JC \\
& B_0 \ar[u] \\
A \ar[r] \ar[uu] & A_0 \ar[u] \ar@{=}[r] & A/J
}
$$
whose vertical arrows are of finite type such that
\begin{enumerate}
\item $\Spec(C) \to \Spec(A)$ is \'etale over $\Spec(A) \setminus V(I)$,
\item $\Spec(B_0) \to \Spec(A_0)$ is \'etale over
$\Spec(A_0) \setminus V(IA_0)$, and
\item $B_0 \to C_0$ is \'etale and induces an isomorphism
$B_0/IB_0 = C_0/IC_0$.
\end{enumerate}
Then we can fill in the diagram above to a commutative diagram
$$
\xymatrix{
C \ar[r] & C/JC \\
B \ar[u] \ar[r] & B_0 \ar[u] \\
A \ar[r] \ar[u] & A/J \ar[u]
}
$$
with $A \to B$ of finite type, $B/JB = B_0$, $B \to C$ \'etale, and
$\Spec(B) \to \Spec(A)$ \'etale over $\Spec(A) \setminus V(I)$.
\end{lemma}

\begin{proof}
Set $X = \Spec(A)$, $X_0 = \Spec(A_0)$, $Y_0 = \Spec(B_0)$,
$Z = \Spec(C)$, $Z_0 = \Spec(C_0)$. Furthermore, denote
$U \subset X$, $U_0 \subset X_0$, $V_0 \subset Y_0$,
$W \subset Z$, $W_0 \subset Z_0$ the complement of the
vanishing set of $I$. Here is a picture to help visualize the
situation:
$$
\xymatrix{
Z \ar[dd] & Z_0 \ar[l] \ar[d] \\
& Y_0 \ar[d] \\
X & X_0 \ar[l]
}
\quad\quad\quad
\xymatrix{
W \ar[dd] & W_0 \ar[l] \ar[d] \\
& V_0 \ar[d] \\
U & U_0 \ar[l]
}
$$
The conditions in the lemma guarantee that
$$
\xymatrix{
W_0 \ar[r] \ar[d] & Z_0 \ar[d] \\
V_0 \ar[r] & Y_0
}
$$
is an elementary distinguished square, see
Derived Categories of Spaces, Definition
\ref{spaces-perfect-definition-elementary-distinguished-square}.
In addition we know that
$W_0 \to U_0$ and $V_0 \to U_0$ are \'etale. The morphism
$X_0 \subset X$ is a finite order thickening as $J$ is assumed nilpotent.
By the topological invariance of the \'etale site
we can find a unique \'etale morphism $V \to X$ of schemes
with $V_0 = V \times_X X_0$ and we can lift the given morphism
$W_0 \to V_0$ to a unique morphism $W \to V$ over $X$.
See \'Etale Morphisms, Theorem
\ref{etale-theorem-remarkable-equivalence}.
Since $W_0 \to V_0$ is separated, the morphism $W \to V$ is separated too,
see for example
More on Morphisms, Lemma \ref{more-morphisms-lemma-deform-property}.
By Pushouts of Spaces, Lemma
\ref{spaces-pushouts-lemma-construct-elementary-distinguished-square}
we can construct an elementary distinguished square
$$
\xymatrix{
W \ar[r] \ar[d] & Z \ar[d] \\
V \ar[r] & Y
}
$$
in the category of algebraic spaces over $X$. Since the base change
of an elementary distinguished square is an elementary distinguished
square (Derived Categories of Spaces, Lemma
\ref{spaces-perfect-lemma-make-more-elementary-distinguished-squares})
we see that
$$
\xymatrix{
W_0 \ar[r] \ar[d] & Z_0 \ar[d] \\
V_0 \ar[r] & Y \times_X X_0
}
$$
is an elementary distinguished square. It follows that there is a
unique isomorphism $Y \times_X X_0 = Y_0$ compatible with the two
squares involving these spaces because
elementary distinguished squares are pushouts (Pushouts of Spaces, Lemma
\ref{spaces-pushouts-lemma-elementary-distinguished-square-pushout}).
It follows that $Y$ is affine by
Limits of Spaces, Proposition \ref{spaces-limits-proposition-affine}.
Write $Y = \Spec(B)$. It is clear that $B$ fits into the desired diagram
and satisfies all the properties required of it.
\end{proof}

\end{document}
